\section{Methodology}
\label{sec:methodology}
\subsection{Thiết lập biến và các nhóm trên bàn cờ}
Với mỗi ô $(i,j)$, tạo một biến Boolean $x_{ij}$, đúng khi và chỉ khi ô đó chứa hậu. ID được cấp theo
\begin{equation}
\operatorname{id}(i,j)=iN+j+1,\qquad 0\le i,j<N.
\label{eq:id}
\end{equation}
Do đó có $N^2$ biến chính. Với $N=4$, bảng ID là
\[
\begin{array}{c|rrrr}
 & j=0&j=1&j=2&j=3\\\hline
i=0&1&2&3&4\\
i=1&5&6&7&8\\
i=2&9&10&11&12\\
i=3&13&14&15&16
\end{array}
\]
Mỗi ID là một biến Boolean, không phải số hậu được đặt.

Định nghĩa các nhóm:
\begin{align}
R_i&=(x_{i0},\ldots,x_{i,N-1}),&
C_j&=(x_{0j},\ldots,x_{N-1,j}),\\
D^-_d&=\{x_{ij}:i-j=d\},&
D^+_s&=\{x_{ij}:i+j=s\}.
\end{align}
Chỉ đưa đường chéo dài từ hai ô trở lên vào tập $\mathcal D$ cần mã hóa. Bỏ đường chéo một ô không mất ràng buộc vì AMO một biến luôn thỏa.

\subsection{Công thức SAT đầy đủ: ALO và AMO}
\subsubsection{Đúng một hậu trên mỗi hàng}
Với từng hàng $i$, thêm ALO:
\begin{equation}
A_R=\bigwedge_{i=0}^{N-1}\left(\bigvee_{j=0}^{N-1}x_{ij}\right).
\label{eq:alo-row}
\end{equation}
Sau đó thêm AMO của nhóm $R_i$ bằng encoding được chọn. ALO buộc hàng có hậu, AMO cấm hai hậu trên cùng hàng; kết hợp hai phần cho đúng một hậu.

\subsubsection{Đúng một hậu trên mỗi cột}
Với từng cột $j$, thêm
\begin{equation}
A_C=\bigwedge_{j=0}^{N-1}\left(\bigvee_{i=0}^{N-1}x_{ij}\right)
\label{eq:alo-col}
\end{equation}
và AMO của $C_j$. Trong đợt thực nghiệm mới, ALO cột bật cho cả năm encoding. Về mặt tập nghiệm, nếu mỗi hàng có đúng một hậu và mỗi cột tối đa một hậu, ALO cột là hệ quả: $N$ hậu ở $N$ cột không trùng nhau phải phủ hết cột. Tuy nhiên giữ các clause này có thể thay đổi suy diễn, nên tùy chọn phải thống nhất khi so sánh.

\subsubsection{Tối đa một hậu trên mọi đường chéo}
Áp dụng AMO cho từng nhóm $D^-_d$ và $D^+_s$ thuộc $\mathcal D$. Không thêm ALO cho đường chéo. Ví dụ các nhóm $(x_{00},x_{11},x_{22},x_{33})$ và $(x_{03},x_{12},x_{21},x_{30})$ là hai đường dài nhất khi $N=4$, nhưng các đường ngắn hơn cũng phải được mã hóa.

\subsubsection{Ghép CNF và tính đúng của mô hình}
Gọi $E_L(L,Y_L)$ là encoding AMO trên nhóm $L$, dùng biến phụ riêng $Y_L$. Công thức toàn bộ là
\begin{equation}
\begin{split}
F_E(X,Y)={}&A_R\land A_C\\
&{}\land\bigwedge_{i=0}^{N-1}E_L(R_i,Y_{R_i})
\land\bigwedge_{j=0}^{N-1}E_L(C_j,Y_{C_j})\\
&{}\land\bigwedge_{D\in\mathcal D}E_L(D,Y_D).
\end{split}
\label{eq:board}
\end{equation}
Mỗi nhóm có bộ biến phụ riêng; các nhóm chỉ chia sẻ biến ô cờ khi cùng chứa ô đó. Một bộ cấp ID bắt đầu từ $N^2+1$ được dùng xuyên suốt để tránh trùng biến phụ.

Nếu $F_E$ thỏa, ALO/AMO hàng cho đúng $N$ hậu, AMO cột và đường chéo loại mọi cặp tấn công. Ngược lại, một bố trí hợp lệ thỏa tất cả ALO/AMO; nhờ~\eqref{eq:existential}, mỗi nhóm có phép gán phụ hợp lệ, và vì các bộ phụ không giao nhau nên ghép được thành phép gán thỏa $F_E$. Như vậy SAT tương ứng một bố trí N-Queens hợp lệ.

\subsection{Khung chung của thuật toán mã hóa}
Trong các phần dưới, $X=(x_1,\ldots,x_k)$ là \emph{một nhóm AMO}, không phải toàn bộ $N^2$ biến. Ký hiệu $a_E(k)$ và $c_E(k)$ là số biến phụ và clause AMO của nhóm. Tất cả encoding dùng
\begin{equation}
\EO_E(X)=\left(\bigvee_{t=1}^{k}x_t\right)\land E(X,Y).
\label{eq:eo-encoding}
\end{equation}
ALO thêm đúng một clause và không thêm biến phụ. Hàng/cột dùng~\eqref{eq:eo-encoding}; đường chéo chỉ dùng $E$. Các nhóm $k\le1$ không cần clause AMO.

\subsection{Binomial / Pairwise encoding}
\subsubsection{AMO}
Với mọi cặp đầu vào khác nhau, thêm
\begin{equation}
E_{\mathrm{pair}}(X)=
\bigwedge_{1\le i<j\le k}(\neg x_i\lor\neg x_j).
\end{equation}
Hai biến cùng đúng làm sai clause của chính cặp đó. Nếu tối đa một biến đúng, mọi cặp đều có ít nhất một biến sai, nên toàn bộ clause thỏa. Không có biến phụ:
\begin{equation}
a_{\mathrm{pair}}(k)=0,\qquad
c_{\mathrm{pair}}(k)=\binom{k}{2}.
\end{equation}
Mã dùng \texttt{combinations(line,2)} để duyệt các cặp không lặp.

\subsubsection{ALO và ví dụ bốn đầu vào}
Exactly-One thêm $(x_1\lor x_2\lor x_3\lor x_4)$ vào
\begin{align*}
&(\neg x_1\lor\neg x_2)\land(\neg x_1\lor\neg x_3)
\land(\neg x_1\lor\neg x_4)\\
{}\land{}&(\neg x_2\lor\neg x_3)\land(\neg x_2\lor\neg x_4)
\land(\neg x_3\lor\neg x_4).
\end{align*}
Nhóm có sáu clause AMO, bảy clause nếu thêm ALO. Mô hình $N=4$ có 16 biến và 84 clause khi bật ALO trên cả hàng và cột.

\subsection{Binary encoding}
\subsubsection{AMO bằng mã nhị phân riêng}
Đặt $b=\lceil\log_2 k\rceil$ và cấp $b$ biến phụ $y_0,\ldots,y_{b-1}$. Biến $x_t$ nhận mã của chỉ số $t-1$. Gọi
\begin{equation}
\ell_{t,h}=
\begin{cases}
y_h,&\text{bit }h\text{ của }t-1\text{ bằng }1,\\
\neg y_h,&\text{bit đó bằng }0.
\end{cases}
\end{equation}
Thêm
\begin{equation}
E_{\mathrm{bin}}(X,Y)=
\bigwedge_{t=1}^{k}\bigwedge_{h=0}^{b-1}
(\neg x_t\lor\ell_{t,h}).
\label{eq:binary}
\end{equation}
Một đầu vào đúng ép toàn bộ bit bằng mã của nó. Hai đầu vào khác nhau có mã khác ít nhất một bit, nên nếu cùng đúng thì một biến phụ phải vừa đúng vừa sai. Khi tất cả đầu vào sai, các bit không bị buộc; khi đúng một đầu vào, chọn mã tương ứng. Do đó thỏa điều kiện~\eqref{eq:existential}.

Số đếm của bản cơ bản hiện tại:
\begin{equation}
a_{\mathrm{bin}}(k)=b,\qquad c_{\mathrm{bin}}(k)=kb.
\end{equation}
Code tính $b$ bằng \texttt{(k-1).bit\_length()}. Các mẫu bit không được dùng khi $k$ không là lũy thừa hai không cần bị cấm để bảo đảm AMO.

\subsubsection{ALO và ví dụ}
Với $k=4$, các điều kiện là
\[
x_1\Rightarrow\neg y_0\land\neg y_1,\quad
x_2\Rightarrow y_0\land\neg y_1,\quad
x_3\Rightarrow\neg y_0\land y_1,\quad
x_4\Rightarrow y_0\land y_1.
\]
Mỗi điều kiện tách thành hai clause, tổng tám. Ví dụ $x_2$ tạo $(\neg x_2\lor y_0)$ và $(\neg x_2\lor\neg y_1)$. Exactly-One thêm ALO trực tiếp $\bigvee_t x_t$, không thêm các hàm ý ngược từ mã bit sang đầu vào. Vì ALO đã chọn một đầu vào, phần AMO đủ để chọn duy nhất một.

Slide trang 8 \citep{amoSlides} có tối ưu cho ba đầu vào bằng cách cho một đầu vào đại diện nhiều mã. Bản hiện tại không áp dụng tối ưu này: $k=3$ có sáu clause AMO, thay vì năm trong ví dụ tối ưu. Không dùng công thức của bản tối ưu để diễn giải số đếm của dữ liệu hiện có.

\subsection{Sequential Counter encoding}
\subsubsection{AMO bằng chuỗi tín hiệu}
Với $k\le3$, dùng Pairwise để tránh tạo chuỗi phụ cho nhóm nhỏ. Với $k\ge4$, cấp $k-1$ biến phụ $s_1,\ldots,s_{k-1}$. Các clause là
\begin{align}
&(\neg x_1\lor s_1),\qquad
(\neg x_k\lor\neg s_{k-1}),\nonumber\\
&(\neg x_i\lor s_i),\quad
(\neg s_{i-1}\lor s_i),\quad
(\neg x_i\lor\neg s_{i-1})
\quad(2\le i\le k-1).
\label{eq:seq}
\end{align}
Một đầu vào đúng truyền tín hiệu về phía sau. Nếu $x_u$ và $x_v$ cùng đúng với $u<v$, $x_u$ ép $s_{v-1}$ đúng còn $x_v$ ép nó sai. Nếu AMO thỏa, có thể chọn $s_i$ bằng OR của tiền tố $(x_1,\ldots,x_i)$ để mở rộng phép gán.

Các clause không ép đầy đủ quan hệ hai chiều giữa $s_i$ và OR tiền tố; đây không phải lỗi, vì chỉ cần tồn tại phép gán phụ. Với $k\ge4$, số đếm của chuỗi Sinz \citep{sinz2005} là:
\begin{equation}
a_{\mathrm{seq}}(k)=k-1,\qquad
c_{\mathrm{seq}}(k)=2+3(k-2)=3k-4.
\end{equation}

\subsubsection{ALO và ví dụ}
Với bốn đầu vào, hai clause biên và sáu clause giữa là
\begin{align*}
&(\neg x_1\lor s_1)\land(\neg x_4\lor\neg s_3)\\
{}\land{}&(\neg x_2\lor s_2)\land(\neg s_1\lor s_2)
\land(\neg x_2\lor\neg s_1)\\
{}\land{}&(\neg x_3\lor s_3)\land(\neg s_2\lor s_3)
\land(\neg x_3\lor\neg s_2).
\end{align*}
Nhóm có ba biến phụ và tám clause AMO. Exactly-One thêm ALO trực tiếp $(x_1\lor x_2\lor x_3\lor x_4)$, tổng chín clause. Phần chuỗi khớp trang 10--11 của slide \citep{amoSlides}. Ngưỡng Pairwise ở $k\le3$ là lựa chọn bổ sung của cài đặt: chi phí lần lượt là $(a,c)=(0,1)$ cho $k=2$ và $(0,3)$ cho $k=3$. Đây không phải khẳng định tối ưu toàn cục; ngay tại $k=4$, Pairwise vẫn có ít clause hơn chuỗi, nhưng cấu hình đang đánh giá đã cố định ngưỡng dừng bằng 3.

\subsection{Commander đệ quy}
\label{sec:commander}
\subsubsection{AMO qua các nhóm đại diện}
Gọi $E_C(X)$ là mã hóa Commander đang được đánh giá. Nếu $k\le3$, mã hóa dừng bằng Pairwise. Nếu $k>3$, chia $X$ theo thứ tự thành $m=\lceil k/3\rceil$ nhóm $G_1,\ldots,G_m$, mỗi nhóm có tối đa ba biến. Cấp một commander $c_j$ cho mỗi nhóm. Ở tầng hiện tại, phát ba loại ràng buộc:
\begin{align}
&\bigwedge_{\{x,z\}\subseteq G_j}(\neg x\lor\neg z),
&&\text{AMO nội nhóm},\label{eq:commander-inner}\\
&\bigwedge_{x\in G_j}(\neg x\lor c_j),
&&\text{đầu vào kéo theo commander},\\
&\left(\neg c_j\lor\bigvee_{x\in G_j}x\right),
&&\text{commander kéo theo ALO nội nhóm}.
\end{align}
Sau đó gọi đệ quy $E_C(c_1,\ldots,c_m)$ để bảo đảm tối đa một nhóm được kích hoạt. Các clause AMO nội nhóm ở~\eqref{eq:commander-inner} có hai literal, không gắn thêm điều kiện $\neg c_j$. Hai chiều liên kết cho $c_j\leftrightarrow\bigvee_{x\in G_j}x$.

Thuật toán luôn kết thúc vì $m<k$ với $k>3$. Bộ cấp biến phụ dùng chung giữa mọi tầng và mọi nhóm AMO trên bàn cờ. Không tái sử dụng ID của một nhóm cho nhóm khác.

\subsubsection{ALO và tính đúng}
Với AMO, không thêm $c_1\lor\cdots\lor c_m$: đường chéo được phép không có hậu. Khi cần Exactly-One trên hàng/cột, thêm duy nhất ALO trực tiếp $x_1\lor\cdots\lor x_k$ lên các đầu vào ban đầu. ALO này cùng liên kết $x\Rightarrow c_j$ tự bảo đảm có commander đúng, nên không cần ALO riêng trên từng tầng.

Chứng minh theo quy nạp trên $k$. Trường hợp dừng Pairwise là đúng. Với $k>3$, hai đầu vào đúng trong cùng nhóm vi phạm AMO nội nhóm; nếu thuộc hai nhóm khác nhau, chúng ép hai commander đúng và vi phạm AMO đệ quy. Ngược lại, nếu không quá một đầu vào đúng, đặt mỗi $c_j$ bằng OR của nhóm tương ứng. Các clause nội nhóm và liên kết đều thỏa, tập commander cũng có không quá một biến đúng. Theo giả thiết quy nạp, có thể mở rộng phép gán cho các tầng cao hơn. Do đó~\eqref{eq:existential} được bảo đảm.

\subsubsection{Ví dụ và số đếm}
Với $k=4$, hai nhóm là $(x_1,x_2,x_3)$ và $(x_4)$. Tầng đầu có ba clause Pairwise, bốn clause $x_t\Rightarrow c_j$, hai clause chiều ngược và một clause $(\neg c_1\lor\neg c_2)$ tại tầng dừng. Tổng cộng hai biến phụ và mười clause AMO; Exactly-One có mười một clause. Nhóm một phần tử vẫn có cả $(\neg x_4\lor c_2)$ và $(\neg c_2\lor x_4)$ theo cài đặt hiện tại.

Ví dụ $k=10$ làm rõ đệ quy nhiều tầng: đầu vào chia thành các nhóm có kích thước $3,3,3,1$, tạo bốn commander; bốn commander lại chia thành nhóm 3 và 1, tạo hai commander mới; tập hai biến cuối dùng Pairwise. Đặt
\[
I(k)=3\lfloor k/3\rfloor+\binom{k\bmod3}{2},\qquad m=\lceil k/3\rceil.
\]
Với $k>3$, công thức là
\begin{align}
a_C(k)&=m+a_C(m),\\
c_C(k)&=I(k)+k+m+c_C(m).
\label{eq:commander-count}
\end{align}
Điều kiện dừng: $a_C(k)=0$, $c_C(k)=\binom k2$ khi $k\le3$. Ví dụ $k=10$ có $a_C(10)=4+2=6$ và $c_C(10)=9+10+4+10=33$.

Slide trang 12--13 \citep{amoSlides} giới thiệu nhóm và commander, với ví dụ dùng Pairwise trên tập commander. Việc lặp lại phép chia nhóm là một mở rộng của nguyên lý đó; họ mã hóa Commander cũng mô tả cách tổ chức đệ quy \citep{klieber2007}. Ngưỡng ba biến, nhóm tối đa ba phần tử và liên kết chiều ngược là các lựa chọn cụ thể cần ghi rõ khi tái lập số đếm.

\subsection{Product đệ quy}
\label{sec:product}
\subsubsection{AMO bằng hai tập tọa độ phụ}
Nếu $k\le4$, dùng Pairwise và không tạo biến phụ. Với $k>4$, đặt
\begin{equation}
p=\lceil\sqrt{k}\rceil,\qquad q=\lceil k/p\rceil.
\end{equation}
Cấp $p$ biến hàng phụ $u_1,\ldots,u_p$ và $q$ biến cột phụ $v_1,\ldots,v_q$. Đây là lưới phụ của \emph{một nhóm AMO}, không phải bàn cờ N-Queens. Đặt mỗi $x_t$ vào một ô riêng:
\[
\alpha(t)=1+\left\lfloor\frac{t-1}{q}\right\rfloor,\qquad
\beta(t)=1+((t-1)\bmod q).
\]
Mỗi đầu vào phát hai clause $(\neg x_t\lor u_{\alpha(t)})$ và $(\neg x_t\lor v_{\beta(t)})$. Cùng mã hóa Product được gọi tiếp trên cả tập $U$ và $V$:
\begin{equation}
\begin{split}
E_P(X)={}&E_P(U)\land E_P(V)\\
&{}\land\bigwedge_{t=1}^{k}
[(\neg x_t\lor u_{\alpha(t)})\land(\neg x_t\lor v_{\beta(t)})].
\end{split}
\label{eq:product}
\end{equation}
Hai tập tọa độ nhỏ hơn đầu vào với $k>4$, nên quá trình dừng tại các nhóm Pairwise. Các ô dư khi $pq>k$ không cần biến đầu vào giả.

\subsubsection{ALO và tính đúng}
Exactly-One tiếp tục dùng ALO trực tiếp $\bigvee_t x_t$. Không thêm ALO trên tọa độ hoặc hàm ý ngược $u_a\land v_b\Rightarrow x_t$. Đặc biệt, không được buộc có tọa độ đúng trong một nhóm AMO đường chéo chỉ vì muốn áp dụng cách mã hóa Exactly-One.

Hai đầu vào khác nhau nằm ở hai ô lưới khác nhau, nên khác ít nhất một tọa độ. Nếu cả hai đúng, chúng ép hai biến hàng hoặc hai biến cột cùng đúng, trái với AMO đệ quy. Ngược lại, nếu không có đầu vào đúng, chọn mọi tọa độ sai; nếu có một đầu vào đúng, chọn đúng hàng và cột tương ứng. Các tập tọa độ thỏa AMO nên theo quy nạp đều có phép gán phụ mở rộng. Cùng với tính đúng của Pairwise ở điểm dừng, điều này chứng minh~\eqref{eq:existential}.

\subsubsection{Ví dụ và số đếm}
Với $k=5$, lưới phụ có $p=3,q=2$. Các đầu vào lần lượt nằm ở $(1,1),(1,2),(2,1),(2,2),(3,1)$. Có mười clause liên kết và $\binom32+\binom22=4$ clause tại hai tập tọa độ. Tổng cộng năm biến phụ, 14 clause AMO và 15 clause nếu thêm ALO.

Với $k=25$, tầng đầu là lưới $5\times5$; mỗi tập năm tọa độ lại dùng lưới $3\times2$ rồi dừng bằng Pairwise. Với $k>4$:
\begin{align}
a_P(k)&=p+q+a_P(p)+a_P(q),\\
c_P(k)&=2k+c_P(p)+c_P(q).
\label{eq:product-count}
\end{align}
Điều kiện dừng: $a_P(k)=0$, $c_P(k)=\binom k2$ khi $k\le4$. Do đó $a_P(25)=10+5+5=20$ và $c_P(25)=50+14+14=78$. Với $k=4$, bản này dùng Pairwise: không có biến phụ và chỉ sáu clause AMO.

Slide trang 14--16 \citep{amoSlides} trình bày phân rã AMO theo hai tập tọa độ và lựa chọn ALO trực tiếp. Ví dụ trên slide sử dụng Pairwise cho tọa độ; cài đặt hiện tại tiếp tục phân rã khi tập tọa độ có hơn bốn biến. Vì thế cần phân biệt nguyên lý Product với lựa chọn một tầng hoặc đệ quy khi so sánh số biến và clause.

\subsection{Công thức đếm cho toàn bàn cờ}
\begin{table}[pos=H]\centering\small
\caption{Số biến phụ và clause của một nhóm AMO. Các ngưỡng áp dụng cho đúng cài đặt thực nghiệm.}
\label{tab:group-counts}
\begin{tabular}{@{}lll@{}}\toprule
Encoding & $a_E(k)$ & $c_E(k)$\\\midrule
Binomial & $0$ & $\binom{k}{2}$\\
Binary & $b=\lceil\log_2 k\rceil$ & $kb$\\
Sequential, $k\le3$ & $0$ & $\binom{k}{2}$\\
Sequential, $k\ge4$ & $k-1$ & $3k-4$\\
Commander, $k\le3$ & $0$ & $\binom{k}{2}$\\
Commander, $k>3$ & $m+a_C(m)$ & $I(k)+k+m+c_C(m)$\\
Product, $k\le4$ & $0$ & $\binom{k}{2}$\\
Product, $k>4$ & $p+q+a_P(p)+a_P(q)$ & $2k+c_P(p)+c_P(q)$\\
\bottomrule\end{tabular}
\end{table}
Các đại lượng $m,I(k),p,q$ được định nghĩa trong hai mục trước. Với $k\le1$, mọi mã hóa có chi phí AMO bằng 0. Exactly-One thêm một clause ALO dài $k$ cho nhóm không rỗng.

Với $N\ge2$, có $2N$ hàng/cột dài $N$ và hai đường chéo dài $N$, tổng $2N+2$ nhóm dài $N$. Với mỗi $2\le k<N$, mỗi hướng chéo có hai đường dài $k$, tổng bốn nhóm. Đặt $d=1$ khi bật ALO cột và $d=0$ khi tắt. Số biến và clause đầy đủ là
\begin{align}
V_E(N)&=N^2+(2N+2)a_E(N)+4\sum_{k=2}^{N-1}a_E(k),\\
C_E(N)&=N(1+d)+(2N+2)c_E(N)+4\sum_{k=2}^{N-1}c_E(k).
\label{eq:whole-count}
\end{align}
Trong toàn bộ đợt thực nghiệm, $d=1$. Hai công thức tính CNF trước tiền xử lý của solver và không tính clause học trong quá trình tìm kiếm.

\begin{table}[pos=H]\centering\small
\caption{Mốc kiểm tra kích thước CNF đầy đủ tại $N=4$, có ALO hàng và cột.}
\begin{tabular}{@{}lrrrr@{}}\toprule
Encoding & Biến chính & Biến phụ & Tổng biến & Clause\\\midrule
Binomial & 16&0&16&84\\
Binary & 16&32&48&120\\
Sequential & 16&30&46&104\\
Commander & 16&20&36&124\\
Product & 16&0&16&84\\
\bottomrule\end{tabular}
\end{table}
Ví dụ Binomial: tám clause ALO, mười nhóm AMO dài 4, bốn nhóm dài 3 và bốn nhóm dài 2 cho $8+10\cdot6+4\cdot3+4\cdot1=84$. Product có cùng kích thước ở $N=4$ vì mọi nhóm đều ở ngưỡng Pairwise. Cùng số biến/clause không đồng nghĩa cùng thứ tự nạp hoặc cùng thời gian chạy.

Từ các công thức, CNF toàn bàn cờ có số clause $\Theta(N^3)$ với Binomial, $\Theta(N^2\log N)$ với Binary và $\Theta(N^2)$ với Sequential, Commander đệ quy, Product đệ quy. Đây là bậc tăng kích thước biểu diễn, không phải độ phức tạp thời gian tìm nghiệm của Glucose3.

\subsection{Quy trình sinh, nạp, giải và kiểm chứng}
\begin{table}[pos=H]\centering\small
\caption{Luồng xử lý SAT của một lượt thực nghiệm hiện tại.}
\label{tab:sat-process}
\begin{tabular}{@{}p{.08\linewidth}p{.85\linewidth}@{}}\toprule
Bước & Xử lý\\\midrule
1 & Tính kích thước CNF đầy đủ bằng công thức. Cấu hình \texttt{max\_clauses=0} không loại trường hợp nào do số clause.\\
2 & Khởi tạo worker, bắt đầu ngân sách toàn lượt, tạo backend SAT và bộ cấp ID phụ từ $N^2+1$.\\
3 & Sinh ALO hàng/cột, AMO hàng/cột/đường chéo; gom tối đa 8192 clause mỗi lô.\\
4 & Ghi số clause đã sinh, nạp lô vào solver, ghi số clause đã nạp và giải phóng list của lô.\\
5 & Khi CNF nạp xong, gọi solver. Bộ giám sát vẫn theo dõi thời gian và RSS của worker cùng các tiến trình con.\\
6 & Nếu có model, lấy các biến ô cờ đúng, chuyển thành vector $q$ và kiểm tra~\eqref{eq:check}.\\
7 & Tiến trình cha nhận kết quả hoặc dừng lượt vượt giới hạn, ghi bản ghi theo định danh lượt chạy.\\
\bottomrule\end{tabular}
\end{table}
Sinh và nạp xen kẽ tránh giữ toàn bộ CNF dưới dạng list Python, nhưng solver vẫn cần bộ nhớ cho công thức và trạng thái tìm kiếm. Với $N=4$, các ID biến chính đúng $3,5,12,14$ biểu diễn nghiệm $(2,0,3,1)$. Biến phụ không được coi là hậu khi giải mã.

Hàm \texttt{estimate\_cnf} tính kích thước CNF hoàn chỉnh theo~\eqref{eq:whole-count}; đó không phải phép ngoại suy thời gian. Hai bộ đếm \texttt{clauses\_generated} và \texttt{clauses\_loaded} phản ánh lượng công việc đã báo về. Nếu timeout trước khi nạp xong, kích thước đầy đủ vẫn giữ nguyên, còn tiến độ có thể nhỏ hơn và chỉ là bản cập nhật gần nhất. Không gộp hai loại đại lượng này vào cùng một cột không chú thích.

Các kiểm tra cần có cho một encoder gồm tính khả thỏa của mọi đầu vào có không quá một biến đúng, tính bất khả thỏa khi ít nhất hai biến đúng, kiểm tra nghiệm bàn cờ, và đối chiếu số clause/literal/biến thực sinh với công thức. Báo cáo này kiểm tra lại toàn bộ vector nghiệm trong dữ liệu và so số clause sinh/nạp của các lượt SAT hoàn tất với công thức. CSV không có phép đếm độc lập toàn bộ ID biến và literal thực sinh trong từng lượt, nên không khẳng định các đại lượng đó đã được đo trực tiếp ở mọi $N$.

\subsection{Mô hình CP và CP-SAT dùng để đối sánh}
Cả CP-SAT và CP Optimizer tạo $N$ biến nguyên
\[
q_i\in\{0,\ldots,N-1\}.
\]
Mỗi biến đã đại diện một hậu ở hàng $i$, nên không cần thêm ALO/AMO hàng dưới dạng clause. Ba ràng buộc là
\begin{equation}
\begin{gathered}
\AD(q_0,\ldots,q_{N-1}),\\
\AD(q_0+0,\ldots,q_{N-1}+N-1),\\
\AD(q_0-0,\ldots,q_{N-1}-(N-1)).
\end{gathered}
\end{equation}
Ràng buộc đầu loại cùng cột; hai ràng buộc sau loại cùng đường chéo. Mô hình có $N$ biến nguyên và ba ràng buộc mức cao trước khi backend phân rã. Hai backend dùng một worker và seed 0 trong đợt thực nghiệm; time limit native lấy ngân sách còn lại. Cấu trúc đầu vào giống nhau không có nghĩa hai backend có thuật toán lan truyền hoặc tìm kiếm giống nhau.

\subsection{Mô hình ILP/MIP dùng để đối sánh}
Cả CPLEX MIP và Gurobi dùng $N^2$ biến nhị phân $z_{ij}$ và
\begin{align}
\sum_{j=0}^{N-1}z_{ij}&=1&&(0\le i<N),\\
\sum_{i=0}^{N-1}z_{ij}&=1&&(0\le j<N),\\
\sum_{i-j=d}z_{ij}&\le1&&(1-N\le d\le N-1),\\
\sum_{i+j=s}z_{ij}&\le1&&(0\le s\le2N-2).
\end{align}
Đẳng thức hàng/cột tương ứng Exactly-One; bất đẳng thức đường chéo tương ứng AMO. Mô hình không có hàm mục tiêu không hằng, nên một nghiệm hợp lệ đã hoàn thành mục tiêu của bài. Code giữ cả các đường chéo một ô, cho $2N+2(2N-1)=6N-2$ ràng buộc tuyến tính trước presolve.

Không gọi các ràng buộc này là clause CNF và không so trực tiếp một biến nguyên miền $N$ giá trị với một biến Boolean. So sánh hiệu năng đánh giá cả mô hình lẫn backend đang dùng. \texttt{native\_vars} và \texttt{native\_constraints} là kích thước giao diện mô hình, không phản ánh cấu trúc nội bộ sau presolve.
